\documentclass[11pt]{article}

\usepackage[T1]{fontenc}
\usepackage{lmodern}
\usepackage[a4paper,margin=29mm]{geometry}
\usepackage{amsmath,amssymb,amsthm,mathtools}
\usepackage{microtype}
\usepackage{enumitem}
\usepackage{xcolor}
\usepackage[colorlinks=true,linkcolor=blue!45!black,citecolor=blue!45!black,
  urlcolor=blue!55!black]{hyperref}

\newtheorem{theorem}{Theorem}[section]
\newtheorem{proposition}[theorem]{Proposition}
\newtheorem{lemma}[theorem]{Lemma}
\newtheorem{corollary}[theorem]{Corollary}
\newtheorem{definition}[theorem]{Definition}
\newtheorem{remark}[theorem]{Remark}

\newcommand{\C}{\mathbb C}
\newcommand{\R}{\mathbb R}
\newcommand{\N}{\mathbb N}
\newcommand{\T}{\mathbb T}
\newcommand{\HH}{\mathbb H}
\newcommand{\cK}{\mathcal K}
\newcommand{\cM}{\mathcal M}
\newcommand{\cP}{\mathcal P}
\newcommand{\cR}{\mathcal R}
\newcommand{\D}{\partial}
\newcommand{\Graph}{\operatorname{Graph}}
\newcommand{\supp}{\operatorname{supp}}
\newcommand{\dtrdeg}{\operatorname{dtrdeg}}
\newcommand{\diag}{\operatorname{diag}}
\newcommand{\Beta}{\operatorname{Beta}}

\title{Cofinal Cholesky Selection for Singular\
Hypertranscendental Herglotz Carriers}
\author{Michael Haufschild \and GPT 5.6 Sol \and Fable 5 \and Grok 4.5}
\date{}

\begin{document}

\maketitle

\begin{abstract}
Fix an indeterminate Hamburger moment sequence and its canonical Nevanlinna boundary simplex.
After Cayley transport to the circle, associate with each representing measure the compatible
inverse-Cholesky matrices of its Toeplitz moment sections.  We prove a sharp selection law.  The
complete matrix tower is a faithful topological coordinate for the carrier, but every finite
prefix level contains positive Hilbert cubes and charts on which maximal differential
algebraic independence of complex translates is residual, even with spectral type fixed.  On the
residual singular-continuous phase, the unrenormalized analytic-polynomial graph is dense in the
full source--carrier product.  The same Cholesky tower is the unique compatible triangular repair;
it yields the unitary orthogonal-polynomial realization and necessarily diverges in both operator
norm and normalized determinant.  Thus finite positive Toeplitz geometry, cofinal carrier
selection, and operator closure are three distinct levels of structure.  The result is a general
constraint theorem for hypertranscendental Herglotz carriers; no arithmetic specialization is used
in its proof.
\end{abstract}

\noindent\textbf{2020 Mathematics Subject Classification.}
30D35, 34M15, 42C05, 44A60, 47A06.

\noindent\textbf{Keywords.}
hypertranscendence, Hamburger moment problem, Herglotz function, Toeplitz matrix,
Cholesky factorization, orthogonal polynomials, linear relation.

\section{Introduction}

Hypertranscendence is usually expressed negatively: a meromorphic function is
hypertranscendental when it satisfies no nonzero algebraic differential equation over the chosen
base field.  For a function selected by positivity or moment constraints, however, the useful
question is geometric.  Which part of the constraint selects the carrier, and which finite data
leave differential freedom?

This paper answers that question for a canonical class of Herglotz carriers.  An indeterminate
Hamburger moment sequence has many positive representing measures with the same complete moment
shadow.  Nevanlinna theory organizes a canonical closed simplex of such measures by probability
laws on the real projective boundary.  Each carrier can also be described, after Cayley transport,
by its Toeplitz moment matrices and their inverse-Cholesky factors.  Those matrices are not an
auxiliary encoding: their rows are the coefficients of the orthonormal polynomials in the
monomial frame.

Our main result identifies a finite/infinite threshold for this Hilbert coordinate.  If
$C_N(\nu)$ denotes the inverse-Cholesky matrix through degree $N$, then
\[
  \nu\longmapsto \bigl(C_0(\nu),C_1(\nu),\ldots\bigr)
\]
is a homeomorphism onto its image.  Nevertheless, fixing any one $C_N$ leaves a positive
Hilbert cube and an infinite-dimensional chart.  On a dense residual subset of that exact level, all finite
families of distinct complex translates of the Cauchy transform are differentially algebraically
independent over $\C(z)$.  The chart can be chosen to preserve spectral type.

The second part determines what occurs at a singular carrier.  The canonical map sending a finite
Hardy coefficient vector to the same analytic polynomial in $L^2(\mu)$ is not merely
nonclosable.  Its graph is dense in the entire product.  Ordinary graph closure therefore retains
no incidence between source and carrier.  The inverse-Cholesky tower repairs the graph before
closure and produces the unitary orthogonal-polynomial representation, but its norm and normalized
volume diverge.  In this sense the complete tower is simultaneously the unique triangular
selector and a noncompact boundary coordinate.

The individual ingredients have classical parents.  Nevanlinna parametrization belongs to moment
theory; Verblunsky coordinates, Bernstein--Szeg\H{o} approximation, and Toeplitz determinant
identities belong to OPUC; the closability boundary is due to Yafaev~\cite{Yafaev2016}; and the
regular--singular decomposition of linear relations is developed in
Hassi--Sebesty\'en--de Snoo--Szafraniec~\cite{HassiEtAl2007}.  The new candidate contribution is
the selection theorem joining these structures to residual full-translation differential
independence inside one exact moment constraint.  We keep this boundary explicit throughout.

The paper is organized as follows.  Section~\ref{sec:boundary} constructs the canonical boundary
simplex.  Section~\ref{sec:frame} identifies finite Cholesky frames with finite Cauchy jets.
Section~\ref{sec:aperture} proves that every finite frame level retains residual maximal
translation hypertranscendence.  Section~\ref{sec:selection} proves complete-tower selection.
Section~\ref{sec:graph} determines the singular graph and its forced repair.  The combined law and
its scope are stated in Sections~\ref{sec:combined} and~\ref{sec:scope}.

\section{The canonical Nevanlinna boundary simplex}\label{sec:boundary}

Let $s=(s_n)_{n\geq0}$ be a normalized indeterminate Hamburger moment sequence and let
\[
  \cM(s)=\left\{\nu\in\cP(\R):\int_{\R}x^n\,d\nu(x)=s_n
  \text{ for every }n\geq0\right\}.
\]
Choose a real Nevanlinna matrix
\[
  W(z)=\begin{pmatrix}A(z)&B(z)\\ C(z)&D(z)\end{pmatrix},
  \qquad A(z)D(z)-B(z)C(z)=1,
\]
in the standard normalization.  Its entries are real entire functions of minimal exponential
type~\cite{Akhiezer1965,PrucknerEtAl2023}.  For
$\tau\in\R\cup\{\infty\}=:\mathbf T$, the constant Nevanlinna parameter gives an
$N$-extremal measure $\nu_\tau\in\cM(s)$ with
\begin{equation}\label{eq:nev-family}
  m_\tau(z):=\int_{\R}\frac{d\nu_\tau(x)}{x-z}
  =\frac{A(z)\tau+B(z)}{C(z)\tau+D(z)}.
\end{equation}
On the regular set of $C$, put
\begin{equation}\label{eq:q}
  q(z)=-\frac{D(z)}{C(z)},\qquad
  a(z)=-\frac{1}{C(z)^2}.
\end{equation}
Direct subtraction gives
\begin{equation}\label{eq:mobius-difference}
  m_\tau(z)=m_\infty(z)+\frac{a(z)}{\tau-q(z)}.
\end{equation}

The supports of distinct $N$-extremal measures are disjoint.  Indeed, a common support point for
two finite parameters would satisfy both $C(x)\tau_1+D(x)=0$ and
$C(x)\tau_2+D(x)=0$, hence $C(x)=D(x)=0$, contrary to
$AD-BC=1$.  The same determinant identity separates the parameter $\infty$.

Define the projective support label
\begin{equation}\label{eq:pi}
  \pi(x)=[-D(x):C(x)]\in\mathbf T.
\end{equation}
It is continuous into the projective line and equals $\tau$ on
$\supp\nu_\tau$.  For $\eta\in\cP(\mathbf T)$ define the barycenter
\begin{equation}\label{eq:beta}
  \Beta(\eta)(E)=\int_{\mathbf T}\nu_\tau(E)\,d\eta(\tau).
\end{equation}
Tonelli's theorem for even moments and domination by the next even moment for odd moments show
that $\Beta(\eta)\in\cM(s)$.

\begin{proposition}[Canonical boundary simplex]\label{prop:beta}
The map $\Beta:\cP(\mathbf T)\to\cM(s)$ is continuous and affine, and
\begin{equation}\label{eq:left-inverse}
  \pi_*\Beta(\eta)=\eta.
\end{equation}
Consequently
\[
  \cK_s:=\Beta(\cP(\mathbf T))
\]
is a closed Bauer simplex canonically affinely homeomorphic to
$\cP(\mathbf T)$.
\end{proposition}

\begin{proof}
The weak continuity of $\tau\mapsto\nu_\tau$ is the classical continuity of the
$N$-extremal Nevanlinna family.  Since $\mathbf T$ is compact, integration against this weakly
continuous kernel proves continuity of $\Beta$.  For every bounded continuous $f$ on
$\mathbf T$, support separation gives
\[
 \int f\,d(\pi_*\Beta(\eta))
 =\int_{\mathbf T}\left(\int_{\R}f(\pi(x))\,d\nu_\tau(x)\right)d\eta(\tau)
 =\int_{\mathbf T}f(\tau)\,d\eta(\tau).
\]
This is \eqref{eq:left-inverse}.  Hence $\Beta$ is injective.  A continuous injection from the
compact space $\cP(\mathbf T)$ into the Hausdorff space $\cM(s)$ is a homeomorphism onto its
closed image.  The extreme points are precisely the measures $\nu_\tau$, so the extreme boundary
is compact.
\end{proof}

The construction does not depend on the chosen normalization.  Changing the Nevanlinna matrix
acts on $\mathbf T$ by a real projective transformation, and the relabeling of $\pi$ is canceled by
the opposite relabeling of $\Beta$.  Thus $\cK_s$ and the retraction
$\Beta\pi_*:\cM(s)\to\cK_s$ are intrinsic.

\begin{proposition}[Spectral transport]\label{prop:spectral-transport}
For every $\eta\in\cP(\mathbf T)$ and every bounded nonnegative $w$ with
$\int w\,d\eta=1$,
\begin{equation}\label{eq:density-transport}
  d\Beta(w\eta)(x)=w(\pi(x))\,d\Beta(\eta)(x).
\end{equation}
In particular, positive bounded reweighting preserves the spectral type of the carrier.  Moreover,
if $\eta$ is atomless and singular with respect to projective Lebesgue measure, then
$\Beta(\eta)$ is singular continuous on $\R$.
\end{proposition}

\begin{proof}
For every bounded Borel function $f$, support labeling gives
\[
 \int f(x)w(\pi(x))\,d\Beta(\eta)(x)
 =\int_{\mathbf T}w(\tau)\left(\int f\,d\nu_\tau\right)d\eta(\tau),
\]
which is \eqref{eq:density-transport}.  If $\eta$ is atomless, then for fixed $x$ only the leaf
$\tau=\pi(x)$ can charge $x$; its singleton parameter has zero $\eta$-mass.  Hence
$\Beta(\eta)$ has no atoms.

If $E\subset\mathbf T$ has projective Lebesgue measure zero and $\eta(E)=1$, then
$\Beta(\eta)$ is supported on $\pi^{-1}(E)$.  On the real line, $\pi$ is a nonconstant real
analytic projective map.  After removing its discrete critical and polar set, the line is a
countable union of intervals on which $\pi$ is a local real-analytic diffeomorphism.  The preimage
of $E$ is therefore Lebesgue-null.  Thus $\Beta(\eta)$ is singular and, by the preceding paragraph,
singular continuous.
\end{proof}

\section{The inverse-Cholesky frame coordinate}\label{sec:frame}

Use the Cayley map
\begin{equation}\label{eq:cayley}
  \kappa(x)=\frac{x-i}{x+i}
\end{equation}
from the compactified real line to $\T$.  For $\nu\in\cK_s$, let
$\mu_\nu=\kappa_*\nu$ and write
\begin{equation}\label{eq:circle-moments}
  c_k(\nu)=\int_{\T}\zeta^k\,d\mu_\nu(\zeta),\qquad k\in\mathbb Z.
\end{equation}
Every member of $\cK_s$ has infinite support, so its Toeplitz Gram matrices
\begin{equation}\label{eq:gram}
  G_N(\nu)=\bigl(c_{j-k}(\nu)\bigr)_{0\leq j,k\leq N}
\end{equation}
are positive definite.

Let $C_N(\nu)$ be the unique lower-triangular matrix with positive diagonal such that
\begin{equation}\label{eq:cholesky}
  C_N(\nu)G_N(\nu)C_N(\nu)^*=I_{N+1}.
\end{equation}
If $v_N=(1,z,\ldots,z^N)^t$, the vector $C_Nv_N$ consists of the first $N+1$
orthonormal polynomials of $\mu_\nu$.  Cholesky uniqueness makes the matrices compatible under
leading-block restriction.

\begin{lemma}[Frame--jet equivalence]\label{lem:frame-jet}
For $\nu_1,\nu_2\in\cK_s$ and $N\geq1$, the following are equivalent:
\begin{enumerate}[label=\textup{(\roman*)},leftmargin=*]
  \item $C_N(\nu_1)=C_N(\nu_2)$;
  \item $G_N(\nu_1)=G_N(\nu_2)$;
  \item $c_k(\nu_1)=c_k(\nu_2)$ for $1\leq k\leq N$;
  \item $m_{\nu_1}^{(r)}(-i)=m_{\nu_2}^{(r)}(-i)$ for $0\leq r<N$.
\end{enumerate}
\end{lemma}

\begin{proof}
Equation~\eqref{eq:cholesky} gives
\begin{equation}\label{eq:gram-inverse}
  G_N(\nu)^{-1}=C_N(\nu)^*C_N(\nu).
\end{equation}
Together with positive-diagonal Cholesky uniqueness this proves (i)$\Leftrightarrow$(ii).
The Toeplitz diagonals and $c_{-k}=\overline{c_k}$ give
(ii)$\Leftrightarrow$(iii).  Finally,
\begin{align}
  \kappa(x)^k
    &=\left(1-\frac{2i}{x+i}\right)^k
      =\sum_{r=0}^k\binom{k}{r}(-2i)^r(x+i)^{-r},
      \label{eq:triangular-cayley}\\
  m_\nu^{(r-1)}(-i)
    &=(r-1)!\int_{\R}(x+i)^{-r}\,d\nu(x),\qquad r\geq1.
      \label{eq:cauchy-derivative}
\end{align}
The change of coordinates from the first $N$ Cauchy integrals to
$(c_1,\ldots,c_N)$ is triangular with nonzero diagonal.  This proves
(iii)$\Leftrightarrow$(iv).
\end{proof}

For $\nu_0\in\cK_s$, define the level-$N$ frame fiber
\begin{equation}\label{eq:frame-fiber}
  \mathfrak F_N(\nu_0)
  =\{\nu\in\cK_s:C_N(\nu)=C_N(\nu_0)\}.
\end{equation}
Lemma~\ref{lem:frame-jet} is the first interaction in the main theorem: a finite canonical Hilbert
frame is exactly a finite regular analytic jet, rather than a merely comparable finite datum.

\section{Finite frame levels retain maximal HT}\label{sec:aperture}

We first specify the differential-algebraic property.  A family $(f_h)_{h\in H}$ of meromorphic
functions is \emph{differentially algebraically independent over $\C(z)$} if the jet family
$(\D^rf_h)_{h\in H,r\geq0}$ is algebraically independent over $\C(z)$.  We call a Cauchy
transform \emph{maximally translation-HT} when
\begin{equation}\label{eq:maxht}
  \bigl(m_\nu(z+h_1),\ldots,m_\nu(z+h_d)\bigr)
\end{equation}
is differentially algebraically independent over $\C(z)$ for every finite tuple of distinct
complex shifts.

Fix $\nu_0=\Beta(\eta_0)$ with $\eta_0$ of infinite support.  For $0<\delta<1$ and $N\geq1$,
let $X_{N,\delta}(\eta_0)$ be the set of real $w\in L^\infty(\eta_0)$ satisfying
\begin{equation}\label{eq:chart}
  \|w-1\|_\infty<\delta,\qquad
  \int w\,d\eta_0=1,\qquad
  \int w(\tau)m_\tau^{(r)}(-i)\,d\eta_0(\tau)
   =m_{\nu_0}^{(r)}(-i)\quad(0\leq r<N).
\end{equation}
The carrier associated with $w$ is
\begin{equation}\label{eq:weighted-carrier}
  \nu_w=\Beta(w\eta_0).
\end{equation}
All weights are strictly positive.  Hence $w\eta_0$ and $\eta_0$ are equivalent, and
Proposition~\ref{prop:spectral-transport} shows that the carrier spectral type is unchanged.

\begin{lemma}[Infinite positive frame aperture]\label{lem:aperture}
The chart $X_{N,\delta}(\eta_0)$ is an open subset of an infinite-dimensional real affine Banach
space.  Every $\nu_w$ belongs to $\mathfrak F_N(\nu_0)$.
\end{lemma}

\begin{proof}
The tangent space at $w=1$ is the common kernel of the real and imaginary parts of finitely many
continuous linear functionals on $L^\infty(\eta_0;\R)$.  Since $\eta_0$ has infinite support,
$L^\infty(\eta_0;\R)$ is infinite-dimensional; a finite-codimensional closed subspace remains
infinite-dimensional.  Its sufficiently small balls preserve positivity.  The last assertion is
Lemma~\ref{lem:frame-jet}.
\end{proof}

\begin{proposition}[Hilbert-cube frame fibers]\label{prop:hilbert-cube}
If $\eta_0$ is atomless, then every $\mathfrak F_N(\nu_0)$ contains a topologically embedded copy
of the Hilbert cube $[-1,1]^{\N}$.  All carriers in this copy have the same complete Hamburger
moments, the same frame through degree $N$, and the same spectral type as $\nu_0$.
\end{proposition}

\begin{proof}
Write the real normalization and jet constraints defining~\eqref{eq:chart} as integration against
$K$ bounded real functions.  Atomlessness lets us choose pairwise disjoint measurable blocks
$B_j$, each split into $K+1$ subsets of positive measure.  The $K+1$ corresponding indicator
functions have a nonzero real linear combination $v_j$, supported on $B_j$, that annihilates all
$K$ constraints.  Normalize $\|v_j\|_\infty=1$.

For $t=(t_j)_{j\geq1}\in[-1,1]^{\N}$, put
\begin{equation}\label{eq:hilbert-cube}
  w_t=1+\frac{\delta}{2}\sum_{j\geq1}2^{-j}t_jv_j.
\end{equation}
The series converges in $L^\infty$, satisfies every constraint, and remains strictly positive.
Disjoint supports make $t\mapsto w_t$ injective.  The factors $2^{-j}$ make it continuous from the
product topology to the uniform norm.  Hence it is a topological embedding because its domain is
compact and the target is Hausdorff.  Composition with $w\mapsto\Beta(w\eta_0)$ remains injective
by Proposition~\ref{prop:beta}.  Lemma~\ref{lem:frame-jet} and
Proposition~\ref{prop:spectral-transport} give the final assertions.
\end{proof}

For completeness, we isolate the transversality input.  It is stated in the exact form used here.

\begin{lemma}[Evaluation codimension]\label{lem:evaluation-codim}
Let $M_1,\ldots,M_t$ be distinct monomials on a finite-dimensional complex jet space, let
$q_0\geq0$, and set $P=t(q_0+1)$.  Choose distinct complex numbers
$z_1,\ldots,z_R$.  If $R>P+q_0$, then the set of packets
$(x_1,\ldots,x_R)$ for which
\begin{equation}\label{eq:evaluation-matrix}
  \mathcal E(x)_{b,(a,d)}=z_b^dM_a(x_b),\qquad
  1\leq b\leq R,\quad 1\leq a\leq t,\quad 0\leq d\leq q_0,
\end{equation}
has column rank smaller than $P$ has complex codimension at least
$R-P-q_0+1$.
\end{lemma}

\begin{proof}
Introduce the incidence set of packets and projective coefficient vectors $[c]\in\mathbb P^{P-1}$
for which
\[
  \sum_{a=1}^t\sum_{d=0}^{q_0}c_{a,d}z_b^dM_a(x_b)=0
  \qquad(1\leq b\leq R).
\]
For fixed nonzero $c$, the polynomial in $x_b$ is identically zero only when every coefficient
polynomial $\sum_dc_{a,d}z_b^d$ vanishes.  Unless all $c_{a,d}$ vanish, this can occur for at most
$q_0$ of the distinct $z_b$.  Each of the remaining at least $R-q_0$ rows imposes one nonzero
polynomial equation in its own variable $x_b$, hence one complex codimension.  The incidence set
therefore has codimension at least $(R-q_0)-(P-1)$ after allowing the projective coefficient
parameter.  Projection to the packet space cannot increase dimension.  This is the asserted
bound.
\end{proof}

\begin{proposition}[Constrained translate transversality]\label{prop:transversality}
Let $J$ be a finite packet of Cauchy jets frozen in a positive boundary chart of the form
\eqref{eq:chart}.  Suppose the boundary law has infinite support.  Then, on the tangent that
annihilates $J$, every finite packet of translated Cauchy jets whose parameter poles are disjoint
from those of $J$ is onto as a real-linear map.  Consequently the set of weights for which
\eqref{eq:maxht} holds for every tuple of distinct shifts is dense residual in the chart.
\end{proposition}

\begin{proof}
We give the two parts because they are where positivity, the uncountable shift space, and
differential coefficients meet.

By \eqref{eq:mobius-difference}, every $z$-jet of $m_\tau(z)$ is a rational function of $\tau$.
Its only parameter pole is $q(z)$, and the highest derivative contributes a nonzero highest-order
principal part whenever $a(z)q'(z)\neq0$.  Reflected conjugates put the poles belonging to real
functionals in the opposite half-plane.

Let $E_J$ be the real tangent annihilating the frozen packet and normalization.  If a real
functional on a target jet packet annihilates its image from $E_J$, then the corresponding real
combination of target rational functions belongs, $\eta_0$-almost everywhere, to the real span of
the constant and the frozen rational functions.  Infinite support upgrades this equality to a
rational identity.  The target and frozen parameter-pole sets are disjoint.  Comparing principal
parts, from the highest order downwards at every target pole, kills every target coefficient.
Thus the target map is onto.

It remains to pass from a fixed tuple of shifts to all tuples.  Fix a differential relation
pattern: finitely many shifts and jet orders, distinct differential monomials
$M_1,\ldots,M_t$, and a degree bound $q_0$ after clearing coefficients in $\C(z)$.  Put
$P=t(q_0+1)$ and choose $R$ so large that
\[
  R-P-q_0+1
\]
exceeds the complex dimension of the shift configuration.  A relation with this pattern supplies
a nonzero null vector for the evaluation matrix~\eqref{eq:evaluation-matrix}.  By
Lemma~\ref{lem:evaluation-codim}, the bad packet set has a finite algebraic stratification whose
every stratum has complex codimension greater than the shift dimension.

The function $q$ has no nonzero translation period.  More intrinsically, if
\[
 D(z+\omega)C(z)=D(z)C(z+\omega)
\]
for some $\omega\ne0$, then any zero of $C$ propagates along a nonconstant arithmetic progression:
the determinant identity prevents $C$ and $D$ from vanishing together.  Such a progression has
linear counting density, contrary to the minimal exponential type of $C$ by Jensen's formula.
Consequently, for distinct shifts $h_i,h_j$, the entire cross-difference
\[
 D(z+h_i)C(z+h_j)-D(z+h_j)C(z+h_i)
\]
is not identically zero.  The derivative numerator
\[
 W(z)=D(z)C'(z)-D'(z)C(z)
\]
is also nonzero as an entire function.  Otherwise $D/C$ would be locally constant on a ball where
$C\ne0$, hence globally projectively constant by the identity theorem, contradicting a zero of
$C$ at which $D\ne0$.

For a finite target-shift packet and finite frozen-pole packet $(t_k)$, take the product of the
shifted denominators $C(z+h_i)$, all cross-differences above, all frozen-avoidance factors
$D(z+h_i)+t_kC(z+h_i)$, and all noncriticality factors $W(z+h_i)$.  This is a nonzero entire
function, so it is nonzero at points of arbitrarily large imaginary part.  At every such row all
target parameter poles are regular, pairwise distinct, disjoint from the frozen poles, and satisfy
$q'\ne0$.  Nonvanishing persists on an open neighborhood of the shift configuration.  We may
therefore cover the distinct-shift configuration space by countably many relatively compact
patches; compactness gives a finite open atlas of such rows on each patch.

The surjectivity proved in the first paragraph supplies, at every shift configuration, a
finite-dimensional perturbation family that is submersive onto all row-jet coordinates.
Surjectivity persists locally; compactness reduces a configuration patch to finitely many such
families.  Parametric transversality~\cite[Chapter~3]{Hirsch1976}, applied successively to the
finitely many smooth algebraic strata, makes the row-evaluation map transverse to the bad set for
a dense set of weights in every neighborhood.  The codimension choice makes its preimage in the
shift patch empty.  Normal convergence of the Cauchy jets on compact upper-half-plane packets
makes this avoidance open.

There are countably many relation patterns and compact configuration patches.  Their open dense
avoidance sets have dense residual intersection by Baire's theorem.  If a relation is first given
as an upper-end germ, the meromorphic identity theorem continues it to the chosen evaluation
rows.  If it has coefficients in $\C(z)$, clearing the finitely many denominators puts it into one
of the enumerated patterns.  Hence the intersection has property~\eqref{eq:maxht}.
\end{proof}

\begin{theorem}[Finite-frame freedom]\label{thm:finite-frame}
For every $N\geq1$, the exact frame fiber $\mathfrak F_N(\nu_0)$ contains the
infinite-dimensional positive chart
$\{\nu_w:w\in X_{N,\delta}(\eta_0)\}$.  Within that chart, maximally translation-HT carriers form
a dense residual set.  Every carrier in the chart has the same complete Hamburger moments, the
same frame through degree $N$, and the same pure-point, singular-continuous, or absolutely
continuous spectral type as $\nu_0$.  If $\eta_0$ is atomless, the frame fiber also contains the
positive Hilbert cube of Proposition~\ref{prop:hilbert-cube}.
\end{theorem}

\begin{proof}
Combine Lemma~\ref{lem:aperture}, Proposition~\ref{prop:hilbert-cube}, and
Proposition~\ref{prop:transversality}.  Exact moment
preservation follows from barycentring measures that all represent $s$.
Proposition~\ref{prop:spectral-transport} preserves the transported spectral type.
\end{proof}

\section{The complete tower selects the carrier}\label{sec:selection}

Define
\begin{equation}\label{eq:theta}
  \Theta:\cK_s\longrightarrow\prod_{N\geq0}\operatorname{Tri}_{N+1}^{+}(\C),
  \qquad \Theta(\nu)=(C_N(\nu))_{N\geq0},
\end{equation}
where $\operatorname{Tri}_{N+1}^{+}(\C)$ denotes lower-triangular matrices with positive diagonal.

\begin{theorem}[Complete-frame selection]\label{thm:selection}
The map $\Theta$ is a homeomorphism of $\cK_s$ onto its image.  Equivalently, for every
$\nu_0\in\cK_s$,
\begin{equation}\label{eq:intersection}
  \bigcap_{N\geq0}\mathfrak F_N(\nu_0)=\{\nu_0\}.
\end{equation}
\end{theorem}

\begin{proof}
If $\Theta(\nu_1)=\Theta(\nu_2)$, Lemma~\ref{lem:frame-jet} gives equality of every positive
circle moment.  Negative moments are their conjugates.  Trigonometric polynomials are uniformly
dense in $C(\T)$, so $\mu_{\nu_1}=\mu_{\nu_2}$.  Cayley inversion yields
$\nu_1=\nu_2$.  This proves injectivity and \eqref{eq:intersection}.

For fixed $N$, $C_N$ depends continuously on the first $N$ circle moments: the moment map is weakly
continuous and positive-diagonal Cholesky factorization is continuous on the positive-definite
cone.  Thus $\Theta$ is continuous in the product topology.  The domain $\cK_s$ is compact and
the target is Hausdorff, so a continuous injection is a homeomorphism onto its image.
\end{proof}

Theorems~\ref{thm:finite-frame} and~\ref{thm:selection} give an exact threshold.  Every finite
projection of the canonical coordinate retains an infinite-dimensional residual-HT modulus,
whereas the coherent projective limit selects one carrier.  No appeal to cardinality or informal
``insufficiency of finite data'' is involved.

\section{Singular graph saturation and the forced repair}\label{sec:graph}

Let $\mu$ be an infinite-support probability measure on $\T$.  On
$\ell^2(\mathbb Z_+)$ define the densely defined polynomial map
\begin{equation}\label{eq:raw-map}
  A_\mu:c_{00}(\mathbb Z_+)\longrightarrow L^2(\mu),
  \qquad A_\mu(a)=\sum_{n\geq0}a_n\zeta^n.
\end{equation}
Its quadratic form has the Toeplitz matrices $G_N(\mu)$ as finite sections.  Yafaev's theorem
states that this form, equivalently $A_\mu$, is closable exactly when
$\mu$ is absolutely continuous with respect to normalized Lebesgue measure~\cite{Yafaev2016}.
For pure singularity one obtains the exact opposite endpoint.

\begin{theorem}[Maximal graph saturation]\label{thm:saturation}
If $\mu\perp m$ is a purely singular probability measure, then
\begin{equation}\label{eq:saturation}
  \overline{\Graph(A_\mu)}
  =\ell^2(\mathbb Z_+)\times L^2(\mu).
\end{equation}
\end{theorem}

\begin{proof}
Suppose $(x,g)$ is orthogonal to $\Graph(A_\mu)$.  Testing the monomials gives, up to the fixed
inner-product convention,
\begin{equation}\label{eq:orthogonal-coefficients}
  x_n=-\int_{\T}\overline{g(\zeta)}\zeta^n\,d\mu(\zeta),
  \qquad n\geq0.
\end{equation}
The complex measure $d\sigma=\overline g\,d\mu$ is finite by Cauchy--Schwarz.  Since
$x\in\ell^2$, one half-line of Fourier coefficients of $\sigma$ belongs to $\ell^2$.
The Riesz brothers theorem implies that $\sigma$ is absolutely continuous; see, for example,
\cite[Chapter~4]{Duren1970} or the form used in~\cite[Theorem~2.6]{Yafaev2016}.  But
$\sigma\ll\mu\perp m$, so $\sigma=0$.  Hence $g=0$ $\mu$-almost everywhere and
$x=0$ by \eqref{eq:orthogonal-coefficients}.  The graph has zero orthogonal complement, which is
equivalent to \eqref{eq:saturation}.
\end{proof}

In the language of linear relations, the closure is maximally multivalued.  The general fact that
the closure of a singular relation is a Cartesian product is classical~\cite{HassiEtAl2007}; the
proof above identifies the full product directly for the present Toeplitz weld.

The same $C_N$ that selected the carrier in Section~\ref{sec:selection} is the unique compatible
triangular repair of this graph.  From \eqref{eq:gram-inverse},
\begin{equation}\label{eq:norm-law}
  \|C_N\|=\lambda_{\min}(G_N)^{-1/2},
  \qquad
  \det C_N=(\det G_N)^{-1/2}.
\end{equation}
Put $D_N=\det G_N$.  The extended Szeg\H{o} limit theorem in its determinant-ratio form
states that, for an infinite-support finite positive measure
$d\mu=w\,dm+d\mu_s$,
\begin{equation}\label{eq:szego-ratio}
  \lim_{N\to\infty}\frac{D_{N+1}}{D_N}
  =\exp\left(\int_{\T}\log w\,dm\right),
\end{equation}
with the value zero when the logarithmic integral is $-\infty$~\cite[Theorem~1]{Lyons2003}.
Equivalently, $w$ is the Radon--Nikod\'ym derivative of the absolutely continuous part of $\mu$.
Our formal interface uses exactly this general target; the specialization to $w=0$ for a singular
measure is part of the machine-checked deduction below, not part of the cited hypothesis.
The following elementary passage is included because it is part of the new geometric deduction,
not part of the cited Szeg\H{o} input.

\begin{lemma}[Ratio-to-root collapse]\label{lem:ratio-root}
If $a_N>0$ and $a_{N+1}/a_N\to0$, then
$a_N^{1/(N+1)}\to0$.
\end{lemma}

\begin{proof}
Fix $0<q<1$.  Eventually $a_{N+1}\leq q a_N$, hence for a suitable $M$ and every $N\geq M$,
$a_N\leq a_Mq^{N-M}$.  Taking $(N+1)$-st roots gives
$\limsup_N a_N^{1/(N+1)}\leq q$.  Letting $q\downarrow0$ proves the claim.
\end{proof}

For pure singularity $w=0$ almost everywhere, so \eqref{eq:szego-ratio} and
Lemma~\ref{lem:ratio-root} give
\begin{equation}\label{eq:szego}
  (\det G_N)^{1/(N+1)}\longrightarrow0.
\end{equation}
Since the geometric mean of the positive eigenvalues dominates their minimum,
$\lambda_{\min}(G_N)\leq(\det G_N)^{1/(N+1)}$.  Therefore
\begin{equation}\label{eq:divergence}
  \lambda_{\min}(G_N)\longrightarrow0,\qquad
  \|C_N\|\longrightarrow\infty,\qquad
  (\det C_N)^{1/(N+1)}\longrightarrow\infty.
\end{equation}

Write $C_N=(c^{(N)}_{jk})_{0\leq j,k\leq N}$ and define
\[
  \varphi_n(\zeta)=\sum_{k=0}^n\overline{c^{(n)}_{nk}}\,\zeta^k.
\]
The leading-principal coherence of the frames says that the polynomial attached to any earlier
row is unchanged at every later level.  Moreover $C_NG_NC_N^*=I$ says exactly that the first
$N+1$ such row polynomials are orthonormal in $L^2(\mu)$.  Since $C_N$ is invertible, those row
polynomials span the same finite-dimensional subspace as $1,\zeta,\ldots,\zeta^N$.  Thus the
stable diagonal rows assemble, without a separate limiting choice, into an isometry
\begin{equation}\label{eq:unitary-map}
  U_\mu:\ell^2(\mathbb Z_+)\longrightarrow L^2(\mu),
  \qquad U_\mu e_n=\varphi_n.
\end{equation}
For pure singularity this map is onto.  If $g$ were orthogonal to all analytic polynomials, the
measure $\overline g\mu$ would have a vanishing Fourier half-line.  F.~and M.~Riesz would make it
absolutely continuous, while it remains singular, so it must vanish.

We have proved the exact three-level geometry
\begin{equation}\label{eq:trichotomy}
\begin{array}{c|c}
\text{finite raw section}&G_N>0\\
\text{unweighted cofinal closure}&\overline{\Graph(A_\mu)}
  =\ell^2\times L^2(\mu)\\
\text{coherently repaired limit}&U_\mu\text{ unitary},\quad \|C_N\|\to\infty.
\end{array}
\end{equation}

For comparison, the Bernstein--Szeg\H{o} measure
\begin{equation}\label{eq:bs}
  d\mu^{[N]}(e^{i\theta})=
  \frac{d\theta}{2\pi|\varphi_N(e^{i\theta})|^2}
\end{equation}
has positive real-analytic density, rational Carath\'eodory transform, and the same first $N$
Verblunsky coefficients as $\mu$~\cite{Simon2005}.  Hence
$G_N(\mu^{[N]})=G_N(\mu)$ and it carries exactly the same finite repair $C_N$.
Every finite model is a bounded, bounded-below weld, while the singular limit has the saturated
raw graph \eqref{eq:saturation}.  What fails is uniform control of the compatible repairs.

\section{The cofinal frame-selection law}\label{sec:combined}

Atomless singular probabilities form a dense $G_\delta$ in $\cP(\mathbf T)$ and are transported
to singular-continuous carriers by Proposition~\ref{prop:spectral-transport}.  The
maximal-translation-HT condition is also residual after transport by $\Beta$.  Indeed, each fixed
relation-pattern avoidance condition on a compact shift patch is open by normal convergence.  It
is dense in every weak neighborhood: first insert an arbitrarily small infinite-support component
into the boundary law, then use a positive $L^\infty$ chart contained in that neighborhood and
Proposition~\ref{prop:transversality}.  Countably many patterns and compact patches give the
residual set by Baire's theorem.  Intersecting the two residual sets gives a residual subset of
$\cK_s$.

\begin{theorem}[Cofinal Cholesky selection]\label{thm:main}
For every indeterminate normalized Hamburger moment sequence, the canonical simplex $\cK_s$
contains a residual set $\cR_s$ of singular-continuous carriers with maximally translation-HT
Cauchy transforms.  For every $\nu_0\in\cR_s$:
\begin{enumerate}[leftmargin=*,label=\textup{(\arabic*)}]
  \item the compatible inverse-Cholesky tower $\Theta(\nu_0)$ determines $\nu_0$ uniquely and
  topologically;
  \item every finite prefix fiber $\mathfrak F_N(\nu_0)$ contains a positive Hilbert cube and a
  chart of infinite dimension with residual maximal translation-HT and fixed singular-continuous
  spectral type;
  \item the unrenormalized analytic-polynomial graph has full-product closure;
  \item the same tower is the unique compatible positive-diagonal triangular repair, produces a
  unitary carrier representation, and obeys the divergence laws~\eqref{eq:divergence};
  \item every Bernstein--Szeg\H{o} model carries the exact finite repair through its truncation
  degree, although all such models are absolutely continuous and differentially algebraic.
\end{enumerate}
\end{theorem}

\begin{proof}
The residual spectral and differential phase was just described.  Item~(1) is
Theorem~\ref{thm:selection}; item~(2) is Theorem~\ref{thm:finite-frame}; items~(3)--(4) follow
from Theorem~\ref{thm:saturation}, Cholesky uniqueness, and
\eqref{eq:divergence}--\eqref{eq:unitary-map}; item~(5) is the
Bernstein--Szeg\H{o} prefix identity following \eqref{eq:bs}.
\end{proof}

The theorem is a positive necessity law.  It does not merely assert that finite Toeplitz data are
insufficient.  It identifies what a faithful carrier construction must retain: a coherent frame
tower, graph norm, boundary relation, or source marking that reconstructs the same unitary
incidence before taking the cofinal closure.  In the singular phase this datum cannot remain
uniformly equivalent to the monomial Hardy frame.

\section{Scope, comparison, and arithmetic specialization}\label{sec:scope}

The statement is about hypertranscendental Herglotz carriers in general.  It uses no primes,
Euler product, functional equation, or zeta zeros.  The differential-algebraic component should be
compared with results on hypertranscendence forced by fixed difference equations, such as
Adamczewski--Dreyfus--Hardouin~\cite{ADH2021}, and with algebraic independence for functions
governed by independent operators~\cite{ADHW2024}.  Here the functions are selected generically
inside an exact positive moment constraint; they are not assumed to solve a difference equation.

The operator statement also has a precise classical boundary.  Graph saturation itself is a
specialized singular-relation phenomenon, not an independent novelty claim.  The OPUC determinant
and approximation formulas are likewise classical.  The candidate new content of
Theorem~\ref{thm:main} is the fact that one and the same canonical matrix tower has all three
roles:
\begin{enumerate}[leftmargin=*]
  \item its finite projections define exact fibers retaining residual maximal HT;
  \item its coherent projective limit is a faithful topological selector;
  \item on the generic singular phase it is the unique divergent triangular repair of a fully
  saturated raw graph.
\end{enumerate}
These roles are not created by independently adjoining two encodings; they are consequences of
the frame--jet equivalence in Lemma~\ref{lem:frame-jet}.

There is a conditional consequence for any Riemann-hypothesis architecture that first normalizes
the logarithmic derivative of the completed zeta function as a Herglotz/Toeplitz carrier.  If that
normalization has a purely singular representing measure, then finite positive Toeplitz sections
and ordinary graph convergence cannot select the source-faithful carrier.  A prime--archimedean
construction must derive the complete compatible frame tower, or an equivalent boundary marking,
before closure.  This is a constraint on such a route, not a proof of positivity and not a proof of
the Riemann hypothesis.

Two limitations are essential.  First, residuality does not provide a canonical explicit member
of $\cR_s$.  Second, the priority of the combined theorem requires specialist comparison across
moment theory, OPUC, linear relations, and differential algebra.  Database searches that fail to
find the exact conjunction do not establish novelty.  These limitations prevent the classical
ingredients from being presented as new and separate the proved theorem from its significance
assessment.

\begin{thebibliography}{99}

\bibitem{ADH2021}
B. Adamczewski, T. Dreyfus, and C. Hardouin,
\emph{Hypertranscendence and linear difference equations},
J. Amer. Math. Soc. \textbf{34} (2021), 475--503.

\bibitem{ADHW2024}
B. Adamczewski, T. Dreyfus, C. Hardouin, and M. Wibmer,
\emph{Algebraic independence and linear difference equations},
J. Eur. Math. Soc. \textbf{26} (2024), 1899--1932.

\bibitem{Akhiezer1965}
N. I. Akhiezer,
\emph{The Classical Moment Problem and Some Related Questions in Analysis},
Oliver \& Boyd, Edinburgh, 1965.

\bibitem{Duren1970}
P. L. Duren,
\emph{Theory of $H^p$ Spaces},
Academic Press, New York, 1970.

\bibitem{HassiEtAl2007}
S. Hassi, Z. Sebesty\'en, H. S. V. de Snoo, and F. H. Szafraniec,
\emph{A canonical decomposition for linear operators and linear relations},
Acta Math. Hungar. \textbf{115} (2007), 281--307;
\href{https://arxiv.org/abs/math/0611045}{arXiv:math/0611045}.

\bibitem{Hirsch1976}
M. W. Hirsch,
\emph{Differential Topology},
Graduate Texts in Mathematics, vol.~33, Springer, New York, 1976.

\bibitem{Lyons2003}
R. Lyons,
\emph{Szeg\H{o} limit theorems},
Geom. Funct. Anal. \textbf{13} (2003), 574--590;
\href{https://arxiv.org/abs/math/0205052}{arXiv:math/0205052}.

\bibitem{PrucknerEtAl2023}
R. Pruckner, J. Reiffenstein, and H. Woracek,
\emph{An upper bound for the Nevanlinna matrix of an indeterminate moment sequence},
\href{https://arxiv.org/abs/2307.10748}{arXiv:2307.10748} (2023).

\bibitem{Simon2005}
B. Simon,
\emph{OPUC on one foot},
Bull. Amer. Math. Soc. \textbf{42} (2005), 431--460;
\href{https://web.ma.utexas.edu/mp_arc/c/05/05-78.pdf}{MP\_ARC 05-78}.

\bibitem{Yafaev2016}
D. R. Yafaev,
\emph{On semibounded Toeplitz operators},
J. Operator Theory \textbf{77} (2017), 205--216;
\href{https://arxiv.org/abs/1603.06229}{arXiv:1603.06229}.

\end{thebibliography}

\end{document}
